\chapter{Introduction}
\label{ch:intro}

Quantitative alpha research scores assets one at a time. A formulaic
alpha reads a single stock's prices and volumes and emits a number; the
cross-section of those numbers, ranked, is the signal
\cite{kakushadze2016alphas}. Yet the assets themselves are anything but
isolated: equities co-move within sectors and supply chains, and
European power prices are literally wired together by cross-border
transmission lines. The premise of this project is that the coupling
structure is information, and its subject is a careful measurement of
how much.

The measurement question splits in two, and the split matters. That a
graph-augmented model outperforms a non-relational baseline is easy to
show and hard to interpret --- the gain may come from the graph, from
the extra parameters, or from information the baseline never saw. This
project therefore fixes the information set (the same island alphas
every model consumes) and separates two claims: what the \emph{graph}
adds, measured by an unlearned uniform-propagation anchor, and what
\emph{learned attention} adds, measured as the gap between a graph
attention network \cite{velickovic2018gat} and that anchor on the same
topology. Three research questions organise the work:

\begin{itemize}
  \item[RQ1] Does learned attention add value over unlearned relational
    propagation, holding inputs and topology fixed?
  \item[RQ2] Does one GAT kernel and one evaluation harness transfer
    across structurally different graphs --- an estimated correlation
    backbone over US equities and a physical interconnector network over
    European bidding zones --- without manufacturing false positives?
  \item[RQ3] On real power-market data, does the physical topology
    improve forecast skill, and where does its value concentrate?
\end{itemize}

A reader from outside quantitative finance should carry one picture
through what follows (\cref{sec:bg-primer} defines the recurring
vocabulary). Every model in this report is trained on the same task:
score today's cross-section of assets so that the scores rank the next
period's price moves, five trading days ahead for equities, one day
ahead for power. Every model in a comparison also receives the same
per-asset inputs: ten signals computed from each stock's own prices and
volumes on the equity track, and each zone's price dynamics plus its
published fundamentals (load, wind and solar forecasts among them) on
the energy track. In particular, the graph neural networks here are not
a vehicle for feeding extra data into price prediction --- whatever
exogenous drivers exist go to the graph-free baselines too. The
relational models' only privilege is topological: permission to
exchange information along edges, estimated return correlations for
equities and the physical cross-border grid for power. Every headline
number in this report is therefore a margin, the value of that
permission with everything else held fixed.

The answers, in brief. On the equity track, learned attention beats the
uniform anchor in thirty of thirty seeded runs, and the best
configuration reaches an out-of-sample (OOS) IC of $0.0148 \pm 0.0043$ and
Sharpe of $1.37 \pm 0.39$ over five seeds; the resulting composite
passes three of four research gates but never beats the single best
island alpha, a failure this report states as plainly as the successes
(\cref{ch:equity}). On the energy track the alpha framing collapsed
honestly: three nested evaluation artifacts each produced spectacular
numbers before the project's own controls caught them, and under honest
returns the strategy loses money (\cref{ch:energy}). Reframed as
forecasting, the same track produced the project's cleanest positives:
the interconnector graph lifts node-level price forecast skill by
$+0.131$ over a no-graph baseline with identical drivers, and on the
irreducibly relational target --- cross-border spreads --- message
passing beats a both-endpoint baseline in five of five seeds
(\cref{sec:en-edge}).

Five contributions follow from this, each defended in a named part of
the report. First, an isolation methodology for relational factors:
matched no-learning anchors, four research gates, and structural
leakage controls, defined once in \cref{ch:methodology} and reused
unchanged by every experiment. Second, the equity result itself ---
attention's margin over naive propagation as a distribution over seeds,
devices, and configurations rather than a point estimate, together with
an interpretability analysis showing the mechanism to be a gentle,
stationary reweighting of mean pooling (\cref{ch:equity}). Third, a
worked cautionary example: the energy alpha arc of
\cref{ch:energy}, in which the project's own controls catch three
artifacts in sequence, ending in a verdict --- rank correlation can be
positive while the strategy loses money --- that generalises beyond
power markets. Fourth, the forecasting reframe and its control-validated
graph value, concentrated on the relational target
(\cref{sec:en-forecast,sec:en-edge}). Fifth, the platform engineering
that makes the rest reproducible: two seams that let island and
relational factors share one pipeline, warehouse marts built from
research runs, and a test suite that pins the leakage-critical logic
(\cref{ch:platform}).

A note on standards. This is a Master's capstone report, and the working
assumption throughout is that its most valuable output is not a
headline Sharpe but a defensible measurement process. Wherever a result
died under scrutiny, the scrutiny and the death are in the text; the
negative results are load-bearing, not confessional.

The report proceeds as follows. \Cref{ch:background} places the work in
the alpha-research, GNN, and electricity-forecasting literatures.
\Cref{ch:platform} describes the platform and the two seams.
\Cref{ch:methodology} fixes the measurement apparatus.
\Cref{ch:equity,ch:energy} report the two tracks.
\Cref{ch:discussion} confronts threats to validity, and
\cref{ch:conclusion} closes with what a successor project should do
differently.
